\section{A second-order Lucas correction: ten positive variables, degree 304}
\label{sec:lucas-q1}

The principal result of this section is a fixed integer polynomial in ten
strictly positive integer variables, with exact degree $304$ and an evaluator
using $87$ binary arithmetic operations.  The arithmetic change repairs a
specific failed direction in the latest inspected submissions: with the
smaller scale $Y=qs$, the elementary first-Pell congruence visibly supplies
only $q^3$, whereas the unknown rank is bounded by $q^4$.  Retaining the next
term of the Lucas polynomial supplies enough precision to recover the index.
The full denominator and the omitted native divisibility are then recovered
in that order.

This is a new step within the inherited native/Pell encoding.  The proof
below composes it with the explicitly retained ordinary calibration of that
encoding to finite syntactic ZFC refutability.  The complete formula, fixed
coefficient bytes and their finite reconstruction are supplied.  No new
end-to-end proof-assistant object or proof of consistency is claimed.

\subsection{The complete polynomial}

The only variables are
\[
 J,F,\alpha,i,T,w,s,n,\eta,\zeta\in\mathbb Z_{>0}.
\]
The four fixed integers $\beta,K_0,M_C,M_F$ are the entire directly indexed
BFS943 tuple of the retained parent.  They satisfy
\[
 \beta\ge31\text{ odd},\qquad K_0>0,\qquad
 0<M_C<\beta<M_F<2\beta.
\]
The largest stored port has $592{,}884$ bits.  These are fixed coefficients,
not additional existential variables.  Their complete finite recipe is
retained separately from the mathematical abbreviations below.

Define the following polynomial expressions:
\begin{gather*}
 h=\beta J,\quad q=h+1,\quad C=q-F-\alpha,\quad X=wq,\quad Y=qs,\\
 k=\eta+\zeta,\quad c=kY+\eta,\quad
 A=Y(X+1)+1,\quad \Delta=A^2-1,\quad H=4A-5,\\
 v=XY^2,\qquad E=1+v(v+1)k^2,\\
 R=[h(q-F)+\alpha](q^2-1)+(M_C+qM_F)J,\qquad
 B=(K_0+w)C-F,\\
 U=ic^2,\quad a=\Delta U^2E,\quad F_0=1+a,\quad G=EF_0,\\
 F_*=1+4aF_0,\quad Q=\Delta(F_*-1),\quad
 b=cT-R,\quad V=F_*b-c,\quad K=X(3c-2X)+2.
\end{gather*}
The new corrected numerator is
\begin{equation}
 e_6=6(k-R-1)-v(R+1)(R-1)(R+3).
 \label{eq:lucas-e6}
\end{equation}
For the principal equation put $L=bV$, and define
\begin{align*}
 D_6&=6q^6,\qquad M=hD_6H,\\
 Z&=hD_6\bigl[H(n-LB)+K\bigr]
       -H\bigl[D_6B+h(e_6-6q^4s)\bigr],\\
 R_G&=M^2G.
\end{align*}
The requested single equation is
\begin{equation}
 \boxed{\;P=Q(Z^2-R_GV^2)+(R_G-Z^2)=0.\;}
 \label{eq:lucas-principal}
\end{equation}
Every named expression is a polynomial abbreviation.  No division, factorial,
binomial coefficient, variable exponent or implicit additional witness occurs
in the equation.  Its literal evaluator is \texttt{equation\_bV.json}.

Replacing $L$ by $V^2$ gives a control equation with degree $410$ in $86$
operations.  It is useful for separating the index correction from the sharper
transport estimate, but is dominated by earlier cost--degree points and is
not claimed as a frontier improvement.

\begin{theorem}[Second-order compact consistency equation]
\label{thm:lucas-q1}
For the complete fixed coefficient tuple just specified, the ordinary
arithmetic proof below, composed with the retained native/refutation
calibration, establishes
\[
 \mathrm{ZFC}\vdash
 \left(\neg\exists\mathbf x\in\mathbb Z_{>0}^{10}\;
 P(\mathbf x)=0\right)\longleftrightarrow\operatorname{Con}(\mathrm{ZFC}).
\]
The principal polynomial has exact degree $304$.  Its supplied evaluator has
$53$ multiplications, $18$ additions and $16$ subtractions, for $87$ binary
operations in total, including squares and multiplication by fixed numerals.
\end{theorem}

For the proof only, write
\[
 t=\frac Bh,\quad W=1+hL,\quad \kappa=\frac KH,\quad
 \sigma=\frac{s}{q^2},\quad \bar e=\frac{e_6}{6},\quad
 \rho=\frac{\bar e}{q^6},\quad z=\frac ZM.
\]
These rational quantities need not initially be integral.  Direct polynomial
identities give
\begin{equation}
 P=(Q-1)Z^2-M^2G(QV^2-1),\qquad
 z=n+\kappa+\sigma-Wt-\rho.
 \label{eq:lucas-carrier}
\end{equation}
In particular the appearance of $1/6$ in the proof does not introduce a
rational-coefficient gate: the circuit computes the integer polynomial $e_6$.

\subsection{Roots and early ranks without assuming source signs}

Positivity gives $q\ge32$, $X,Y\ge q$, $k\ge2$, $c>kY$, $A\ge3$,
$U\ge c^2$, $F_*>2c$, $Q>1$ and $M>0$.  For every integer $b$,
$V=F_*b-c$ is nonzero and $|V|>1$.  The exact identities
\[
 F_*=(1+2a)^2,\qquad Q=4\Delta^2U^2G
\]
show that $G$ is coprime to both $Q-1$ and $QV^2-1$.  At a zero of $P$,
the first identity in \eqref{eq:lucas-carrier} yields
$(Q-1)z^2=G(QV^2-1)$.  Every prime valuation of $G$ is even, so
$G=g^2$ with $g>0$ integral.  Since $\gcd(E,F_0)=1$, its positive factors
are separately squares:
\[
 E=\tau^2,\qquad F_0=f_0^2,\qquad g=\tau f_0.
\]
The positive rational $y=|z|/g$ satisfies
\begin{equation}
 QV^2-(Q-1)y^2=1.
 \label{eq:lucas-adjacent}
\end{equation}

The adjacent-Pell descent forces $y$ integral.  To recall its short proof,
write $y=r/m$ in lowest terms.  Then $m^2\mid Q-1$.  If $u=|V|>1$, then
$u<y<Qu/(Q-1)$, and the norm-preserving transformation
\[
 u'=(2Q-1)u-2(Q-1)y,\qquad
 y'=(2Q-1)y-2Qu
\]
has a nonzero integer first coordinate with $|u'|<u$ and preserves the
reduced denominator $m$ of its second coordinate.  Reflection and descent
reach $u=y=1$, forcing $m=1$.  Thus $z=\varepsilon gy$ is integral,
where $\varepsilon\in\{-1,1\}$.

Set $f=1+2a$ and $I=2Ug$.  Pell doubling gives
\begin{equation}
 f^2-\Delta I^2=1,\qquad Q=(\Delta I)^2,\qquad c^2\mid I,
 \qquad V=f^2(cT-R)-c.
 \label{eq:lucas-strong}
\end{equation}
The retained signed all-index folding theorem applies to
\eqref{eq:lucas-adjacent}--\eqref{eq:lucas-strong}, $c\le I$ and
$V\equiv-c\pmod f$.  It requires no sign for $V$ or $R$ and gives
$c=\psi_A(p)$ for a positive odd $p$.  Here $\psi_A(0)=0$,
$\psi_A(1)=1$ and $\psi_A(j+2)=2A\psi_A(j+1)-\psi_A(j)$.
The first norm gives
\[
 k=2\psi_{2v+1}(n_0),\qquad n_0\ge1.
\]
If $n_0=1$, then $k=2$, $\eta=\zeta=1$, and $c=2Y+1$ lies strictly
between the first two main ordinates $1$ and $2A$.  This is impossible.
Consequently
\begin{equation}
 n_0\ge2,\quad k\ge8v+4,\quad
 c>kY>8XY^3\ge8q^4,\quad c>X,\quad K>0.
 \label{eq:lucas-early-growth}
\end{equation}
No sign of $B,C,R,b,V$ or $\bar e$ has been assumed.

\subsection{Absolute transport at the smaller scale}

If $b=0$, equation \eqref{eq:lucas-adjacent} would give
$y^2-c^2=(c^2-1)/(Q-1)$ strictly between zero and one.  Hence $b\ne0$,
$\operatorname{sgn}V=\operatorname{sgn}b$, and $|V|>|b|$.  With
$D_0=\Delta U^2$, one has
\[
 \frac{a^2}{g^2}=\frac{D_0a}{1+a}>\frac{D_0}{2}>9.
\]
Thus $a>3g$, $|V|>3g$ and $y<3|V|/2$.  In particular
$gy<V^2/2$ for the control.

For $L=bV$, the third main index must also be excluded.  Suppose
$c=4A^2-1$.  If $n_0\ge3$, then $k>32v^2$, while $A<2XY$ gives
\[
 c<16X^2Y^2<32X^2Y^5<kY<c,
\]
a contradiction.  Thus $n_0=2$ and $k=8XY^2+4$.  Put
$u=X+2-2Y$.  Direct expansion gives
\[
 \eta=4Y^2(Xu+1)+8Y(X+1)+3-4Y.
\]
For $u\le-1$ this is negative, using $Y(X-1)>2(X+1)$.  For $u\ge2$,
subtracting $k$ leaves at least $4Y^2+8Y(X+1)-4Y-1>0$.  Both contradict
$0<\eta<k$.  The remaining $u=0,1$ would make $2Y-X$ equal to $2$ or $1$,
contradicting $q\mid X,Y$ and $q\ge32$.  Therefore $p\ge5$ and
$c>c_3:=4\Delta+3$.

Put $S=\Delta I$, so $S^2=Q$.  Pell classification of
\eqref{eq:lucas-adjacent} gives $|V|=\chi_S(j)/S$ for a positive odd $j$.
Index one gives $|V|=1$ and is excluded.  At index three,
$|V|=4Q-3$ would give one of
\[
 F_*(b-4\Delta)=c-c_3,\qquad F_*(b+4\Delta)=c+c_3.
\]
Both right sides lie strictly between zero and $2c<F_*$, impossible.
Hence $j\ge5$, $|V|\ge16Q^2-20Q+5>8Q^2$, and
\[
 |b|>\frac{8Q^2-c}{F_*}>4F_*>4g.
\]
It follows that $L=bV=|b||V|>0$ and $gy<3L/8<L/2$.  For either transport
choice we have established
\begin{equation}
 L>0,\qquad gy<L/2,\qquad |b|/q^6<L/2.
 \label{eq:lucas-transport-bounds}
\end{equation}

\subsection{The corrected numerator and all possible rank signs}

The new correction needs a fresh sign argument: its cubic changes sign at
negative ranks.  Set $e_0=k-R-1$ and
$B_R=(R+1)(R-1)(R+3)$.  Since $b-e_0=cT-k+1>0$,
\[
 \rho<\frac b{q^6}-\frac{vB_R}{6q^6}.
\]
When $B_R\ge0$, this implies $\rho<|b|/q^6$.  If $B\le-1$, then
$Wt<-L$ and \eqref{eq:lucas-carrier}, with positive $\kappa,\sigma$ and
\eqref{eq:lucas-transport-bounds}, gives $n<0$.

The integer exceptions are exactly $R\le-4$ and $R=0$.  Write
$P_{\rm mask}=(M_C+qM_F)J>0$.  The exact outer identity is
\begin{equation}
 R=(hC+q\alpha)(q^2-1)+P_{\rm mask}.
 \label{eq:lucas-signed-outer}
\end{equation}
For $R=-r\le-4$, it implies $C<0$ and
\[
 |C|>\frac r{h(q^2-1)}>\frac r{q^3}.
\]
Thus $B=-(K_0+w)|C|-F\le-1-w|C|$ and
$Wt<-(1+w|C|)L$.  The adverse correction is bounded by
\begin{equation}
 \delta_0=\frac{v(r-1)(r+1)(r-3)}{6q^6}
          <\frac{ws^2r^3}{6q^3},
 \label{eq:lucas-adverse}
\end{equation}
because $v=wq^3s^2$.  Here $b=cT+r>r,c$, so
$V>(F_*-1)b$ and either choice of $L$ satisfies
$L>(F_*-1)r^2$.  Also $\Delta>Y^2\ge s^2$ and
$F_*-1=4a(1+a)>s^2$.  Consequently $w|C|L>\delta_0$.  Substitution gives
\[
 n<L/2-(1+w|C|)L+|b|/q^6+\delta_0<0.
\]

For $R=0$, \eqref{eq:lucas-signed-outer} again gives $C<0$, while
$\delta_0=v/(2q^6)=ws^2/(2q^3)$.  Now $b=cT\ge c$ and
$V\ge(F_*-1)b$, so $L>F_*-1$ and
$w|C|L>w(F_*-1)>\delta_0$.  The same contradiction follows.
This handles every integer rank, including the cases where the new cubic
correction has the unfavorable sign.

Therefore every positive zero has $B\ge0$, which forces $C>0$ because
$F>0$.  The retained mask bound now applies:
\begin{equation}
 3q<R<q^4-2.
 \label{eq:lucas-rank-box}
\end{equation}
Equation \eqref{eq:lucas-early-growth} gives $c>2R$, whence $b,V>0$.
The retained normalized rank and signed auxiliary-index theorems have all
their original hypotheses: $c^2\mid I$, $c>2R,c>2p$,
$V\equiv-c\pmod f$ and $V\equiv-R\pmod c$.  They give $p=R$.
In particular $R$ is positive odd; the unchanged mask parity gives even $q$
and odd $J$.

\subsection{The second-order index lemma}

The new argument is elementary and works at every prime, including the primes
dividing the displayed denominator $6$.

\begin{lemma}[Corrected index recovery]
\label{lem:lucas-corrected-index}
At every positive zero of \eqref{eq:lucas-principal},
\[
 n_0=(R+1)/2,\qquad q^6\mid\bar e,
 \qquad \rho\in\mathbb Z_{\ge0},\qquad q^2\mid s.
\]
The conclusions hold as well for the square-transport control.
\end{lemma}

\begin{proof}
The Lucas polynomial has the finite expansion
\begin{equation}
 \psi_{1+2v}(m)=\sum_{j=0}^{m-1}(4v)^j
                    \binom{m+j}{2j+1}.
 \label{eq:lucas-binomial-expansion}
\end{equation}
For a direct derivation, its recurrence gives the formal generating function
\[
 \sum_{m\ge1}\psi_{1+2v}(m)t^{m-1}
 =\frac1{(1-t)^2-4vt}
 =\sum_{j\ge0}\frac{(4v)^jt^j}{(1-t)^{2j+2}}.
\]
Coefficient comparison yields \eqref{eq:lucas-binomial-expansion}.  In
particular, for an explicitly determined nonnegative integer $Z_m$,
\begin{equation}
 2\psi_{1+2v}(m)
 =2m+8v\binom{m+1}{3}+32v^2Z_m.
 \label{eq:lucas-second-order}
\end{equation}
The formula also covers $m=1,2$, when $Z_m=0$.

Put $r_0=(R+1)/2$, an integer by the preceding rank argument.  Then
$B_R=48\binom{r_0+1}{3}$, so $\bar e=e_6/6$ is an integer and
\begin{equation}
 \bar e=k-2r_0-8v\binom{r_0+1}{3}.
 \label{eq:lucas-ebar}
\end{equation}
The integrality of $z$ and the carrier give
$6q^6\mid hH(e_6-6q^4s)$.  Cancelling the literal common factor $6$,
now that \eqref{eq:lucas-ebar} is integral, gives
$q^6\mid hH(\bar e-q^4s)$.  Because $h=q-1$ and
$H=4Y(X+1)-1\equiv-1\pmod q$, one has $\gcd(q,hH)=1$.  Hence
\begin{equation}
 q^6\mid\bar e-q^4s,\qquad q^4\mid\bar e.
 \label{eq:lucas-carrier-congruence}
\end{equation}
No inverse of $6$ modulo $q$ has been taken.

Set $d=2n_0-R-1=2(n_0-r_0)$.  Reducing
\eqref{eq:lucas-second-order} and \eqref{eq:lucas-ebar} modulo $4v$
gives $\bar e\equiv d\pmod{4v}$.  Since $q^3\mid v$, equation
\eqref{eq:lucas-carrier-congruence} first gives $q^3\mid d$.
The exact Taylor difference is
\begin{align}
 6(\bar e-d)
 &=vd\bigl[(2n_0)^2+(2n_0)(R+1)+(R+1)^2-4\bigr]
       +192v^2Z_{n_0}.
 \label{eq:lucas-taylor-upgrade}
\end{align}
Both terms are divisible by $q^6$.  If a prime $\ell\mid q$ has
$a=v_\ell(q)\ge1$, then $v_\ell(6)\le1$, so
\[
 v_\ell(\bar e-d)\ge6a-1\ge4a.
\]
Thus $q^4\mid\bar e-d$, uniformly at $\ell=2,3$ as well.  Together
with \eqref{eq:lucas-carrier-congruence}, this gives $q^4\mid d$.

The first base $2XY^2+1$ exceeds $A$.  Were $n_0\ge R$, monotonicity
in the positive base and index would give $k/2\ge\psi_A(R)=c$,
contrary to $c>k$.  Therefore $1\le n_0<R$, and
\eqref{eq:lucas-rank-box} gives $|d|<R<q^4$.  Hence $d=0$.

Only now return to the full expansion: equations
\eqref{eq:lucas-second-order}--\eqref{eq:lucas-ebar} give
$\bar e=32v^2Z_{r_0}$.  Because $q^3\mid v$, this proves
$q^6\mid\bar e$ and $\rho\ge0$ integral.  Finally,
\eqref{eq:lucas-carrier-congruence} gives $q^6\mid q^4s$, hence
$q^2\mid s$.
\end{proof}

The two congruence steps are essential.  The first establishes divisibility
of the index error by $q^3$; this makes the correction error small in the
relevant prime valuations and upgrades it to the $q^4$ modulus needed for
the ordinary size bound.  The full $q^6$ denominator is recovered after the
index has been identified.

\subsection{Projection and restoration of the native source}

Lemma~\ref{lem:lucas-corrected-index} makes $\sigma=s/q^2$ integral and
restores $q^3\mid Y=qs$ before projection.  The carrier now gives
$\kappa-WB/h\in\mathbb Z$, so $H\mid hK$.
Write $R=2r+1$.  By \eqref{eq:lucas-rank-box}, $r\ge48$.
The retained ratio theorem, uniformly valid at scale one, gives
\[
 Y>X^r/3,\qquad X^{r+1}<3Y(X+1)<H.
\]
The universal Pell projector and the retained entire-product estimate give
\[
 H\mid h(X-2^R)(2^RX+1),\qquad
 |h(X-2^R)(2^RX+1)|<H.
\]
The product is zero, so $X=2^R$ and $\kappa$ is integral.  No possibly
nonunit factor of the product has been cancelled modulo $H$.  The remaining
carrier gives $h\mid WB$, and $W\equiv1\pmod h$ gives $h\mid B$.
Thus the native quotient $t=B/h$ is a nonnegative integer.

All inherited fixed-minus parity and exact half-binomial rounding hypotheses
have now been restored.  The fresh reconstruction uses the original source
base $Y(X+1)+2$, which is distinct from the auxiliary base $A$ used here.
It produces an actual zero of the calibrated native source.  The whole fixed
coefficient tuple's retained native/refutation theorem therefore yields a
finite syntactic ZFC refutation from every positive zero of
\eqref{eq:lucas-principal}.

\subsection{Construction of positive witnesses}

Conversely, a finite refutation supplies the parent's canonical outer data:
$q^3\mid Y$, $X=2^R$, $R\equiv3\pmod4$, $R>3q$ and integral
$t=B/h\ge0$, with the exact half-binomial interval.  Define the new positive
input $s=Y/q$.  Then $\sigma=s/q^2=Y/q^3$ is an integer and at most $Y$.
Using the scale-one bases, put
\begin{gather*}
 c=\psi_A(R),\qquad
 k=2\psi_{2XY^2+1}\bigl((R+1)/2\bigr),\\
 \tau=\chi_{2XY^2+1}\bigl((R+1)/2\bigr),\qquad
 \eta=c-kY,\quad\zeta=k-\eta.
\end{gather*}
The inherited uniform ratio estimate gives $\eta,\zeta>0$.  The full
Lucas expansion gives $\rho\ge0$ integral; projection gives positive
integral $\kappa<Xc$.

A positive multiple of a finite Pell matrix period modulo $c^2\tau$
supplies one seed pair $(f_0,I_0)$ with $c^2\tau\mid I_0$.
Set $i=I_0/(c^2\tau)>0$ and double this pair exactly as in the parent.
It yields \eqref{eq:lucas-strong}, $G=(\tau f_0)^2$ and the canonical
positive auxiliary pair $V,y$, with the positive integral unshifted input
\[
 T=\frac{V+c+Rf^2}{cf^2}.
\]
No private strong pair is reselected after its auxiliary congruences are
specified.  Moreover
\[
 f_0>(A-1)I_0>2Xc,\qquad \kappa<Xc,\qquad \sigma\le Y<Xc,
\]
so $g=\tau f_0>\kappa+\sigma$.  Now set
\[
 n=gy+(1+hL)t+\rho-\kappa-\sigma.
\]
This is a strictly positive integer, and the resulting carrier is $z=gy$.
The polynomial therefore vanishes.  Every one of the ten independent
inputs has been constructed positively.  This completes both existential
directions in Theorem~\ref{thm:lucas-q1}.

\subsection{Exact degree and validation}

The relevant global degrees are
\[
 \begin{array}{c|rrrrrrrr}
 \text{expression}&Y&c&A&\Delta&v&E&U&a\\\hline
 \deg&2&3&4&8&6&14&7&36
 \end{array}
 \qquad
 \begin{array}{c|rrrrrrr}
 \text{expression}&F_*&Q&G&b&V&M&Z\ (L=bV)\\\hline
 \deg&72&80&50&4&76&11&93
 \end{array}
\]
The correction $v(R+1)(R-1)(R+3)$ has degree $18$, below the leading
transport part.  For $L=bV$, the term $-QM^2GV^2$ uniquely dominates the
final expression and has degree $304$.  For $L=V^2$, $\deg Z=165$ and
$QZ^2$ uniquely supplies degree $410$.

At the smaller scale, $cT$ and $R$ both have degree four; the leading part
of $b=cT-R$ must retain both.  Put
\[
 q_0=\beta J,\qquad k_0=\eta+\zeta,\qquad C_0=q_0-F-\alpha,
 \qquad B_0=k_0sT-q_0^2(q_0-F).
\]
The entire highest homogeneous polynomials are
\begin{align*}
 [P_{bV}]_{304}
   &=-9\cdot2^{12}q_0^{128}w^{34}s^{78}i^{14}k_0^{44}B_0^2,\\
 [P_{V^2}]_{410}
   &=9\cdot2^{16}q_0^{166}w^{46}s^{104}i^{20}k_0^{60}B_0^4C_0^2.
\end{align*}
For example, their respective coefficients at
\[
 J^{128}w^{34}s^{80}i^{14}\eta^{46}T^2,\qquad
 J^{166}w^{46}s^{108}i^{20}\eta^{64}T^4F^2
\]
are $-9\cdot2^{12}\beta^{128}$ and
$9\cdot2^{16}\beta^{166}$, both nonzero at the authentic fixed tuple.
These give global degree certificates, rather than sampled line degrees.

The independent exact checker binds all actual evaluator rows through $37$
free-polynomial cuts per circuit.  It recounts the live gates, checks every
fixed coefficient byte, and propagates the complete highest homogeneous
forms: $274$ terms for the
principal equation and $3{,}750$ for the control.  Actual operator corruption
checks reject $175$ modified principal artifacts and $173$ control artifacts.
Normal and optimized Python runs agree.  A separate retained reconstruction
script regenerates all four fixed numerals from their finite recipe and
checks them against the stored bytes.  The Taylor identity
\eqref{eq:lucas-taylor-upgrade} is checked exactly; separate finite Lucas
recurrence controls are explicitly diagnostic.

Those finite checks certify the equations and algebraic certificates.  They
do not replace the unbounded arithmetic proof above or the named inherited
calibration, and do not decide the existence of a ZFC refutation.  The new
work has also received a separate ordinary mathematical review within this
AI-assisted workflow; that is not an outside human referee's endorsement.
